\section{Audit, claim boundary, and the next Gram gate}
\label{sec:tpc58-claims-next}

The point of the preceding construction is to locate, rather than conceal,
the two comparisons that are still missing.  With the final sign convention
of \cref{prop:tpc58-exact-gauge-alignment},
\begin{equation}
 z_H=-\sum_p b,
 \qquad
 h=C_{\mu,H}z_H=-\sum_pv_p,
 \qquad
 A=U_\mu z_H,
 \qquad
 C_1A=h.
 \label{eq:tpc58-final-sign-ledger}
\end{equation}
Thus the physical affine column is \(v_0+h\).  The first missing constant
controls the erasure \(z_H\mapsto h\), and the second controls the sum
\((v_0,h)\mapsto v_0+h\).  They have independent kernels, so they must not
be combined into one unnamed ``Gram loss.''

\subsection{Six-point physical audit}

The native-affine bridge passes the following bookkeeping audit.  Passing
the audit certifies that the operators act on the intended physical object;
it does not certify positivity of either restricted singular value.
\begin{enumerate}[label=\textup{(A\arabic*)},leftmargin=3.2em]
 \item \emph{Literal coefficients.}
       The terminal amplitude is the inherited physical coefficient
       \(x_u=c_u\zeta_{m(u)}\), including its output factor, source shape,
       masks, profiles, cofactor sign, and terminal multiplier.  In the
       TPC--56--57 common gauge, \(x_u\) already contains the factor
       \(\mu(s(u))\).  Equation \eqref{eq:tpc58-final-sign-ledger} therefore
       uses the unsigned second grouping \(C_1A\); no second M\"obius sign is
       inserted.
 \item \emph{Fixed physical shift.}
       Every terminal atom satisfies
       \(p_ur_u-m(u)j_u=h_0\) for the one declared nonzero physical shift.
       The identities in
       \cref{prop:tpc58-exact-gauge-alignment,prop:tpc58-native-exchange-identity}
       are pointwise at this \(h_0\).  No intercept average, frozen-shift
       family, or projective reassembly is substituted for the physical
       column.
 \item \emph{Physical atomic normalization.}
       The terminal diagonal energy, the first grouped residual energy, and
       the native high energy use the inherited raw counting-measure atoms.
       Repeated cofactor labels remain distinct terminal atoms until the
       declared grouping, and the terminal multiplier is not moved into a
       renormalized input norm.  Consequently the two grouping maps compare
       actual physical energies rather than support cardinalities.
 \item \emph{Canonical representation and minimality scope.}
       The first grouping uses the full residual key \(q(u)=(i,s(u))\); the
       second grouping erases only \(s\).  The restricted space
       \(E_H=\Ran T_H\) and the feasible affine pair space
       \(\cF_{\rm aff}\) are the canonical images of one declared coefficient
       map.  Moore--Penrose reductions are restricted to the active Gram
       range, so null coefficient directions are not counted as physical
       generalized eigenvectors.
 \item \emph{Actual active support.}
       The baseline carrier used to define \(T_H\) is fixed independently of
       the row coefficient vector.  Only factors that vanish identically on
       that baseline are removed.  Coefficient-dependent zeros may shrink a
       distinguished support, but they do not redefine \(E_H\), its rank, or
       the uniform constants.  The high residual identification also uses
       the same \(p>y\) cutoff on both the Walsh and terminal sides.  At the
       endpoint this cutoff forces the smaller cofactor range
       \eqref{eq:tpc58-compatible-cofactor-ceiling}; occupancy of that range
       is not inferred from the old \(R\asymp J\) block.
 \item \emph{Endpoint budget.}
       The two new losses are recorded separately as
       \(\lambda_H\) and \(\lambda_{\rm aff}\).  Together with the separately
       recorded compatible-band relocation, selection, cofactor-block,
       terminal, singleton, boundary, and tail losses, they enter the exact
       ledger
       \eqref{eq:tpc58-total-endpoint-exponent}.  The safe condition is the
       strict inequality
       \(\Lambda_{\rm tot}<1/400\); equality is not accepted without a
       constant-level comparison and quantified remainder.
\end{enumerate}

\subsection{Claim table}

\begin{center}
\small
\begin{tabular}{>{\raggedright\arraybackslash}p{0.29\textwidth}
                >{\raggedright\arraybackslash}p{0.60\textwidth}}
\toprule
Statement & Status \\
\midrule
Signed residual-label erasure spectrum and bright/dark decomposition
& Exact finite-dimensional theorem on the declared ambient fibers (L0). \\
Support-only native lower bound
& Sharply false whenever one native output contains at least two active
  residual labels (L0 no-go). \\
Restricted native constant \(\alpha_H\)
& For nonzero \(E_H=\Ran T_H\), an exact least-eigenvalue and kernel
  characterization, with one-way rank/dimension kill tests (L0).
  Positivity and endpoint size are not proved. \\
Affine constant \(\delta_{\rm aff}\)
& Exact restricted singular-value and anti-diagonal kernel
  characterizations on nonzero \(\cF_{\rm aff}\) (L0).  Positivity and
  endpoint size are not proved. \\
Two-stage affine transfer
& Exact implication
  \(\cN_{\rm aff}\ge\delta_{\rm aff}\alpha_H\|z_H\|^2\)
  for vectors in the declared feasible images (L0). \\
Fixed-\(h_0\) terminal/Walsh/native identification
& Exact coordinatewise interface under the inherited injection, cutoff, and
  common-carrier hypotheses (L1). \\
TPC--56 adaptive selected sector
& Valid for its selected pre-terminal object, but it does not transport to
  the undeleted residual or native vector (L0 nontransport theorem). \\
TPC--57 terminal singleton bridge
& Controls only the first grouping \(p\mapsto(i,s)\), conditionally on its
  own occupancy hypotheses and after a compatible smaller cofactor band has
  been activated; it does not control the second grouping
  \((i,s)\mapsto i\) (L0/L1 boundary). \\
Native signed exchange identity
& Exact expansion over unequal residual products at one fixed \(h_0\)
  (L1).  The preceding same-residual factor-exchange estimate does not apply
  to these pairs. \\
Uniform native high lower square
& On nonzero \(\Ran R_H\), equivalent to the negative-spectrum bound
  \(\lambda_{\min}(\Theta_H)>-1\); not proved. \\
Full native affine lower square
& Requires both a native high lower square and independent base--high
  noncancellation; not proved. \\
Endpoint \(1/400\) comparison
& Conditional exponent propagation only.  None of the missing positive
  constants is entered as zero-cost. \\
Parity improvement or twin-prime consequence
& Not claimed. \\
\bottomrule
\end{tabular}
\end{center}

The level terminology is deliberately conservative.  L0 comprises exact
finite-dimensional identities, spectra, restricted minimax laws, and sharp
counterexamples.  L1 is the exact specialization of those maps to the
literal fixed-\(h_0\) carrier.  An L2 advance would require a quantitative
arithmetic estimate proving a sufficiently strong lower bound for the
distinguished physical vector or for the whole canonical coefficient image,
together with every upstream and downstream transfer in
\eqref{eq:tpc58-total-endpoint-exponent}.  No theorem in this paper is
classified as L2.

\subsection{The next gates and their kill tests}

The remaining work begins with one support door and then separates into two
spectral doors.  Compatible-band activation must be tested before either
spectral constant: a Gram theorem on an empty inherited block carries no
endpoint energy.  After activation, the native door precedes the affine
door because an affine estimate cannot reconstruct a high component already
lost in residual-label erasure.

\begin{enumerate}[label=\textbf{Door \Alph*.},leftmargin=3.8em]
 \item \emph{Compatible-band relocation and occupancy.}
       With \(y=X^{267/400+\eta}\), replace the incompatible old block
       \(R\asymp J\) by dyadic blocks contained in
       \[
        r\le X^{133/400-\eta+o(1)}.
       \]
       Prove a lower transfer from the inherited parent diagonal into these
       blocks, retain one fixed \(h_0\), and then verify terminal activation
       and the hypotheses of the TPC--57 undeleted singleton bridge.  The
       relocation loss \(\lambda_{\rm reloc}\) and subsequent block loss
       \(\lambda_{\rm blk}\) remain distinct.

       The immediate kill test is asymptotic emptiness or a best possible
       parent-to-band exponent already exhausting the \(1/400\) slack.  In
       either case the endpoint route stops before any Gram computation.

 \item \emph{Literal restricted-Gram assembly.}
       Assemble the pair
       \[
        R_H=T_H^*T_H,
        \qquad
        N_H=T_H^*C_{\mu,H}^*C_{\mu,H}T_H
       \]
       from the actual fixed-carrier coefficients, after quotienting only
       \(\Ker T_H\).  The first target is a certified census of the negative
       spectrum of \(\Theta_H\), not merely a floating-point plot.  Exact
       sparse elimination, interval eigenvalue bounds, and blockwise
       Schur-complement certificates may be used, provided every block keeps
       the literal normalization and active support.

       The immediate kill test is
       \[
        E_H\cap\Ker C_{\mu,H}\ne\{0\}.
       \]
       If it holds on the actual image, then \(\alpha_H=0\) and the
       coefficient-uniform native door is closed.  A positive numerical
       eigenvalue at finitely many scales is scouting evidence only; a proof
       requires a scale-uniform lower certificate with an explicit loss
       exponent.

 \item \emph{Distinguished signed exchange.}
       If the uniform image contains a dark vector, retain the one inherited
       coefficient vector and estimate the weakest exact target
       \begin{equation}
        2\Re\sum_i\sum_{s<t}A_{i,s}\overline{A_{i,t}}
        \ge -(1-\delta_H)\sum_{i,s}|A_{i,s}|^2
        \label{eq:tpc58-next-distinguished-exchange}
       \end{equation}
       with \(\delta_H\ge X^{-\lambda_H+o(1)}\).  Unlike the TPC--57
       same-residual collision system, the pairs here satisfy
       \(s(u)\ne s(v)\).  Any arithmetic proof must therefore expose a new
       unequal-residual-product relation rather than reusing singleton
       occupancy under a renamed key.

       This door is closed for the endpoint if literal examples force
       \(\delta_H\) to decay faster than the residual endpoint slack.  Such a
       result would be a genuine negative theorem about this route, not a
       claim of cancellation for unrelated coefficient families.

 \item \emph{Base--high affine transversality.}
       Conditional on surviving Door B or C, assemble the feasible pair image
       \(\cF_{\rm aff}\) and test
       \begin{equation}
        \cF_{\rm aff}\cap\{(g,-g):g\in\cH_{\rm out}\}=\{0\}.
        \label{eq:tpc58-next-affine-kill-test}
       \end{equation}
       A nonzero intersection proves \(\delta_{\rm aff}=0\) for the uniform
       route.  Otherwise the target is a quantitative principal-angle lower
       bound
       \(\delta_{\rm aff}\ge X^{-\lambda_{\rm aff}+o(1)}\).  This analysis
       must retain the joint dependence of \(v_0\) and \(h\) on the same
       coefficient vector; treating them as independent ambient vectors can
       only recover the sharp no-go of
       \cref{thm:tpc58-two-independent-no-go}.

 \item \emph{Distinguished affine ratio.}
       If uniform affine transversality fails but the physical coefficient
       vector avoids the anti-diagonal, study the exact pointwise ratio
       \eqref{eq:tpc58-distinguished-delta-aff}.  As at the native gate, a
       distinguished estimate is weaker than a coefficient-uniform theorem
       and must be labeled accordingly.  Its loss exponent still enters the
       same strict endpoint ledger.

 \item \emph{Budget stop and route closure.}
       At each scale record
       \(\lambda_{\rm Gram}=\lambda_H+\lambda_{\rm aff}\) separately from
       the five upstream and two downstream losses.  If the best rigorous
       lower exponents force
       \(\Lambda_{\rm tot}\ge1/400\), the current endpoint route is not
       certified.  At equality, stop unless all constants and remainders are
       closed.  If either exact kernel test is triggered, record the
       corresponding negative result and return to a different physical
       carrier or decomposition; do not replace the failed undeleted vector
       by the adaptive TPC--56 sector.
\end{enumerate}

The most economical next paper is therefore a compatible-band census at the
literal fixed shift: determine whether the smaller \(r\)-range forced by
\(p>y\) carries a quantitatively useful part of the parent energy.  Only on
a surviving nonzero carrier should the subsequent paper construct \(R_H\),
\(N_H\), and the affine pair Gram and perform exact rank and kernel tests.
This order permits a positive theorem or a sharp route-closure theorem
without performing spectral analysis on an empty face.

\paragraph{Conclusion.}
The missing native bridge is now an exact signed alias problem.  TPC--57
can, subject to its own unresolved arithmetic hypotheses, deliver energy
after the first grouping, but residual-label erasure introduces the new
unequal-residual exchange form.  Its optimal restricted constant is
\(\alpha_H\).  Even if that gate is positive, completion of the physical
column encounters the independent affine constant
\(\delta_{\rm aff}\).  Both constants have exact spectral descriptions and
sharp independent zero mechanisms; neither is proved positive on the
endpoint carrier.  The paper therefore closes an L1 interface and exposes
one upstream activation gate and the next two auditable spectral gates,
while making no parity improvement and no twin-prime claim.
